% BOOK-FILE-ID: f-c8997aff583f
% subsection: 連立型

\subsection{連立型}\label{節：連立型}
ここまでは1つの数列について,隣り合う項の関係を考えてきた.
では,2つの数列が互いに絡み合って変化する場合はどうなるだろうか.
一般的な形は
\[
  \begin{cases}
    a_{n+1}=pa_n+qb_n\\
    b_{n+1}=ra_n+sb_n
  \end{cases}
\]
で与えられる.
$a_{n+1}$を求めるには$b_n$が,\,$b_{n+1}$を求めるには$a_n$が必要なので,どちらか一方だけを取り出して基本パターンに当てはめることはできない.
そこで,2つの数列を組み合わせて,等比型になる1つの数列を作ることを目標にしよう.

\subsubsection{和と差で分ける}
まずは係数がきれいに入れ替わっている場合から考える.
例えば
\[
  \begin{cases}
    a_{n+1}=3a_n+b_n\\
    b_{n+1}=a_n+3b_n
  \end{cases}
\]
を眺めてみると,\,$a_n$と$b_n$の立場が入れ替わっているだけである.
こういう対称な式を見たら,まず足したり引いたりしてみるのが常套である.
実際に足すと,
\[
  a_{n+1}+b_{n+1}=4a_n+4b_n=4(a_n+b_n)
\]
となり,和$a_n+b_n$が公比$4$の等比型になる.
引くと
\[
  a_{n+1}-b_{n+1}=2a_n-2b_n=2(a_n-b_n)
\]
となり,差$a_n-b_n$も公比$2$の等比型になる.
和と差が分かれば,あとは$a_n$と$b_n$についての連立方程式を解くだけである.

\begin{bookpoint}
  \textbf{POINT}\par
  係数が対称なら,2式を足し引きして和と差の等比型を作り,連立方程式を解いて元に戻す.
\end{bookpoint}

\noindent \bookblocklink{example-system-1}{problem}{answer}{【例題】}\par\nobreak
\begin{enumerate}[label=(\arabic*),parsep=3pt,format=\bookitemlink{example-system-1}{problem}{answer}]
  \item $a_1=1,\,b_1=1.\ \begin{cases}a_{n+1}=3a_n+b_n\\ b_{n+1}=a_n+3b_n\end{cases}$
\end{enumerate}

\noindent \bookblocklink{example-system-1}{hint}{problem}{【方針】}\par\nobreak
\begin{enumerate}[label=(\arabic*),parsep=3pt,format=\bookitemlink{example-system-1}{hint}{problem}]
  \item 係数$3,1$が入れ替わるので,和と差がそれぞれ等比型になります.和と差の初項を求めてから,\,$a_n,b_n$へ戻します.
\end{enumerate}

\noindent \bookblocklink{example-system-1}{answer}{problem}{【解答】}\par\nobreak
\begin{enumerate}[label=(\arabic*),parsep=3pt,format=\bookitemlink{example-system-1}{answer}{problem}]
  \item $a_n=4^{n-1},\qquad b_n=a_n=4^{n-1}$
\end{enumerate}

\noindent \bookblocklink{example-system-1}{solution}{problem}{【解法】}\par\nobreak
\begin{enumerate}[label=(\arabic*),parsep=3pt,format=\bookitemlink{example-system-1}{solution}{problem}]
  \item 二つの式を足すと$a_{n+1}+b_{n+1}=4(a_n+b_n)$なので,
    \[
      a_n+b_n=(a_1+b_1)4^{n-1}=2\cdot4^{n-1}.
      \booktag{sim-example-sum}{1}
    \]
    引くと$a_{n+1}-b_{n+1}=2(a_n-b_n)$なので,
    \[
      a_n-b_n=(a_1-b_1)2^{n-1}=0.
      \booktag{sim-example-difference}{2}
    \]
    $\bookeqbackref{sim-example-sum}{(1)}+\bookeqbackref{sim-example-difference}{(2)}$から$2a_n=2\cdot4^{n-1}$,
    $\bookeqref{sim-example-sum}{(1)}-\bookeqref{sim-example-difference}{(2)}$から$2b_n=2\cdot4^{n-1}$を得る.
    従って,\,$a_n=b_n=4^{n-1}$である.
\end{enumerate}

\paragraphbf{係数が対称な場合の一般形}
一般に,係数の位置が入れ替わった\bookterm{recurrence}{漸化式}
\[
  \begin{cases}
    a_{n+1}=pa_n+qb_n\\
    b_{n+1}=qa_n+pb_n
  \end{cases}
\]
では,足すと両方の係数が$p+q$,引くと$p-q$とその符号反転になる.
したがって,和と差がそれぞれ等比型になる.\footnote{以下の累乗の式は$p+q,\,p-q$が0でない場合に使う.0になる方は,初項を別に扱い,\,$n\ge2$で0とする.}
\begin{align*}
  &a_{n+1}+b_{n+1}=(p+q)(a_n+b_n),\\
  &a_n+b_n=(a_1+b_1)(p+q)^{n-1}
  \booktag{sim-general-sum}{1}
\end{align*}
\begin{align*}
  &a_{n+1}-b_{n+1}=(p-q)(a_n-b_n),\\
  &a_n-b_n=(a_1-b_1)(p-q)^{n-1}
  \booktag{sim-general-difference}{2}
\end{align*}
$\bookeqbackref{sim-general-sum}{(1)}+\bookeqbackref{sim-general-difference}{(2)}$から,
\[
  \begin{aligned}
    a_n&=\frac12\bigl\{(a_1+b_1)(p+q)^{n-1}\\
    &\qquad+(a_1-b_1)(p-q)^{n-1}\bigr\}
  \end{aligned}
\]
を得る.\,$\bookeqref{sim-general-sum}{(1)}-\bookeqref{sim-general-difference}{(2)}$からは,
\[
  \begin{aligned}
    b_n&=\frac12\bigl\{(a_1+b_1)(p+q)^{n-1}\\
    &\qquad-(a_1-b_1)(p-q)^{n-1}\bigr\}
  \end{aligned}
\]
となる.

\subsubsection{重みを付けて分ける}
では,係数が対称でない一般の連立漸化式ではどうだろうか.
単に足すと,
\[
  a_{n+1}+b_{n+1}=(p+r)a_n+(q+s)b_n
\]
となり,\,$p+r=q+s$でなければ右辺が$a_n+b_n$の定数倍にならない.
足すだけ,引くだけでうまくいかないなら,足し方に重みを付けてやればよい.
つまり,それぞれに係数$\alpha,\beta$を掛けて足した$\alpha a_n+\beta b_n$を考え,これが等比型になるように$\alpha,\beta$を決めるのだ.
\[
  \begin{aligned}
    \alpha a_{n+1}+\beta b_{n+1}
    &=(p\alpha+r\beta)a_n\\
    &\quad+(q\alpha+s\beta)b_n
  \end{aligned}
\]
これを$\gamma(\alpha a_n+\beta b_n)$にそろえるには,係数を比較して
\[
  \begin{cases}
    (p-\gamma)\alpha+r\beta=0\\
    q\alpha+(s-\gamma)\beta=0
  \end{cases}
\]
を満たす重みを選べばよい.\,$\alpha=\beta=0$では何も分からないので,少なくとも一方が0でない組を探す.

まず$\gamma$を決めよう.第1式を$s-\gamma$倍し,第2式の$r$倍を引けば$\beta$が消える.
同様に$\alpha$を消去すると,
\begin{align*}
  \{(p-\gamma)(s-\gamma)-qr\}\alpha&=0,\\
  \{(p-\gamma)(s-\gamma)-qr\}\beta&=0
\end{align*}
を得る.重みは同時に0ではないので,共通の係数が0,すなわち
\[
  \gamma^2-(p+s)\gamma+(ps-qr)=0
\]
が必要になる.これを連立漸化式の\bookterm{characteristic}{特性方程式}と呼ぶ.
形式的には,係数を並べた$\begin{pmatrix}p&q\\r&s\end{pmatrix}$から作られる方程式である.\footnote{\chapterref{章：解法の背景と広がり}で扱う行列の固有値を求める方程式そのものである.}

その解$\gamma_1,\gamma_2$を順に重みの条件へ入れ,非零の組$(\alpha,\beta),(\alpha',\beta')$を選ぶ.
\begin{align*}
  c_n&=\alpha a_n+\beta b_n,\\
  d_n&=\alpha'a_n+\beta'b_n
\end{align*}
と置けば,\,$c_{n+1}=\gamma_1c_n,\ d_{n+1}=\gamma_2d_n$なので,倍率が0でなければ,
\[
  c_n=c_1\gamma_1^{n-1},\qquad
  d_n=d_1\gamma_2^{n-1}
\]
となる.初項は元の$a_1,b_1$から求める.\footnote{倍率が0の場合は,その量の初項を別に扱い,\,$n\ge2$で0とする.}

元の項へは,
\[
  \begin{cases}
    c_n=\alpha a_n+\beta b_n\\
    d_n=\alpha'a_n+\beta'b_n
  \end{cases}
\]
を解いて戻す.一意に復元するには,重みの二組が比例しないこと,すなわち$\alpha\beta'-\alpha'\beta\ne0$が必要である.
\textbf{異なる二根なら,重みは比例しない.}
もし比例すれば,同じ非零の重みが両方の係数条件を満たす.条件を引くと
\[
  (\gamma_1-\gamma_2)\alpha=(\gamma_1-\gamma_2)\beta=0
\]
となり,\,$\gamma_1\ne\gamma_2$より$\alpha=\beta=0$である.これは0でない重みを選んだことに矛盾する.

異なる根は十分条件であり,必要条件ではない.
重解では比例しない二組を取れるかで分かれる.\footnote{例えば$a_{n+1}=2a_n,\ b_{n+1}=2b_n$では$(1,0),(0,1)$を選べる.
  二組を取れない場合は,一方の量を先に求め,残る式へ代入する.\bookproblemref{practice-basic}{25}{標準問25}では,その式を指数の付加項として解く.\booknavref{節：付加項がない場合}{隣接3項型の重解}とも比べられる.}
このように重みを選べば,元の連立漸化式から等比型の数列を2つ作れる.\footnote{解が虚数の場合は\booknavref{節：特性根が虚数の場合}{特性根が虚数の場合}を参照.}

% スマホ版では直前の脚注とPOINTの縦方向のはみ出しを避ける.
\BookMathLayout{}{\pagebreak}
\begin{bookpoint}
  \textbf{POINT}\par
  $\alpha a_n+\beta b_n$が等比型になるように係数比較し,特性方程式$\gamma^2-(p+s)\gamma+(ps-qr)=0$の2解ごとに重みを決める.
\end{bookpoint}

\noindent \hyperlink{book:example-system-1:answer:2}{\textbf{【例題】}}\par\nobreak
\begin{enumerate}[start=2,label=(\arabic*),parsep=3pt,format=\bookitemlink{example-system-1}{problem}{answer}]
  \item $a_1=1,\,b_1=0,\ \begin{cases}a_{n+1}=3a_n+b_n\\ b_{n+1}=2a_n+2b_n\end{cases}$
\end{enumerate}

\noindent \hyperlink{book:example-system-1:problem:2}{\textbf{【方針】}}\par\nobreak
\begin{enumerate}[start=2,label=(\arabic*),parsep=3pt,format=\bookitemlink{example-system-1}{hint}{problem}]
  \item $\gamma$を求めてから重みを選ぶと,\,$a_n-b_n$と$2a_n+b_n$がそれぞれ等比型になります.
\end{enumerate}

\noindent \hyperlink{book:example-system-1:problem:2}{\textbf{【解答】}}\par\nobreak
\begin{enumerate}[start=2,label=(\arabic*),parsep=3pt,format=\bookitemlink{example-system-1}{answer}{problem}]
  \item $a_n=\frac{2\cdot4^{n-1}+1}{3},\qquad b_n=\frac{2\cdot4^{n-1}-2}{3}$
\end{enumerate}

\noindent \hyperlink{book:example-system-1:problem:2}{\textbf{【解法】}}\par\nobreak
\begin{enumerate}[start=2,label=(\arabic*),parsep=3pt,format=\bookitemlink{example-system-1}{solution}{problem}]
  \item $a_1=1,\,b_1=0,\ \begin{cases}a_{n+1}=3a_n+b_n\\ b_{n+1}=2a_n+2b_n\end{cases}$\\
    重みの条件は,
    \[
      \begin{cases}
        (3-\gamma)\alpha+2\beta=0\\
        \alpha+(2-\gamma)\beta=0
      \end{cases}
    \]
    となる.0でない重みを取るための特性方程式は,
    \begin{align*}
      0&=(3-\gamma)(2-\gamma)-2\\
      &=\gamma^2-5\gamma+4=(\gamma-1)(\gamma-4)
    \end{align*}
    なので,\,$\gamma_1=1,\ \gamma_2=4$である.

    $\gamma=1$では$\alpha+\beta=0$となり,\,$(\alpha,\beta)=(1,-1)$を選べる.
    $c_n=a_n-b_n$と置くと,
    \[
      c_{n+1}=c_n,\quad c_1=1
      \quad\Longrightarrow\quad c_n=1
    \]
    である.\,$\gamma=4$では$\alpha=2\beta$となり,\,$(2,1)$を選べる.
    $d_n=2a_n+b_n$と置くと,
    \[
      \begin{aligned}
        &d_{n+1}=4d_n,\quad d_1=2,\\
        &\therefore\quad d_n=2\cdot4^{n-1}
      \end{aligned}
    \]
    を得る.復元は$c_n+d_n=3a_n,\ b_n=a_n-c_n$から,
    \begin{align*}
      a_n&=\frac{2\cdot4^{n-1}+1}{3},\\
      b_n&=\frac{2\cdot4^{n-1}-2}{3}
    \end{align*}
    となる.
\end{enumerate}

\paragraphbf{和と差は重みの一例}\label{解釈：和と差は重みの一例}
\bookproblemref{example-system-1}{1}{最初の例題}へ戻ろう.
$p=s=3,\ q=r=1$を特性方程式へ入れると,
\[
  \gamma^2-6\gamma+8=(\gamma-2)(\gamma-4)=0
\]
となる.\,$\gamma=2$なら$\alpha+\beta=0$より重み$(1,-1)$,
$\gamma=4$なら$\alpha=\beta$より$(1,1)$を選べる.
従って取り出す量は,
\[
  c_n=a_n-b_n,\qquad d_n=a_n+b_n
\]
であり,更新は$c_{n+1}=2c_n,\ d_{n+1}=4d_n$となる.
初項$c_1=0,\ d_1=2$も最初の解法と一致する.
対称な係数なら足し引きが早いが,それは重み$(1,1),(1,-1)$がたまたま見えていたからなのだ.
和と差も,重みを付ける方法の一例なのである.

\paragraphbf{別解：隣接3項型に帰着させる}
ここで扱う別解は,計算が楽になるものではない.
しかしながら,連立型と隣接3項型が同じものであることを示すため,解説する.

$b_n$を消去してみよう.1つ目の式から$qb_n=a_{n+1}-pa_n$である.
1つ目の式の添字を1つ進めて2つ目の式を代入すると,
\begin{align*}
  a_{n+2}&=pa_{n+1}+qb_{n+1}\\
  &=pa_{n+1}+q(ra_n+sb_n)\\
  &=pa_{n+1}+qra_n+s(a_{n+1}-pa_n)\\
  &=(p+s)a_{n+1}-(ps-qr)a_n
\end{align*}
となり,\,$a_n$だけの隣接3項型になるのだ.
その特性方程式
\[
  x^2-(p+s)x+(ps-qr)=0
\]
は,重みを決めるときに現れた特性方程式と全く同じである.
2つの方法は,同じ方程式を別の方向から見ていたのである.
逆に,隣接3項型$a_{n+2}+pa_{n+1}+qa_n=0$も,\,$b_n=a_{n+1}$と置けば$a_n,b_n$の連立型として書ける.
この見方は\chapterref{章：解法の背景と広がり}の行列でもう一度扱う.
